% Copyright 2010-2012  by tgoerlach < tobias.goerlach@uni-hohenheim.de >
% This is a example file for beamerthemehohenheim
% https://bitbucket.org/tobig/hohenheimbeamertheme/
%
% This file may be distributed and/or modified
%
% 1. under the LaTeX Project Public License and/or
% 2. under the GNU Public License.
%
% See the file doc/licenses/LICENSE for more details.


\documentclass[xcolor={table}]{beamer}

\usepackage[ngerman]{babel}
\usepackage[latin1]{inputenc}
\usepackage[T1]{fontenc}
\usepackage{lmodern}


% Document specific stuff ()
\usepackage{etex}% in some cases that helps% eg if you get those no room for newdimen errors
% here we need it for pgfplots
\usepackage{booktabs, url, calc}
\usepackage{phaistos}% Just for the bee
\newcommand{\bee}{\scalebox{0.5}{\PHbee}\ }
\usepackage{tabularx} % tabu is even better
\usepackage{amsmath, amssymb} % ok no commend needed here
\usepackage{listings}% for code % note that beamer wants you to put code only into fragile frames (ref to the beamerusersguide)
\lstdefinestyle{tex}{language=tex, basicstyle={\footnotesize\ttfamily} , breaklines=true}
\usepackage{pgfplots}% best function plotter for tex
 \pgfplotsset{compat=newest}
\pgfplotsset{tick label style={font=\small\sffamily}, label style={font=\small\sffamily}, legend style={font=\footnotesize\sffamily},
}







%-- Theme Stuff
\usepackage[framelist front={o\ }, framelist back={\ }]{hobete}
%--that all what you have to do. 


%\HohenheimLogoKlein{/usr/local/texlive/2011/texmf-dist/doc/generic/pstricks/images/tiger.pdf} %Use the eps provided by the University
%\HohenheimLogoLang{/usr/local/texlive/2011/texmf-dist/doc/generic/pstricks/images/tiger.pdf} % Lang indicates that here mus be insert the "Wort und Bildmarke"
%Klein indicates, that here should only the Bildmarke be inserted. 
%-- You can gain them from the University CI Center
%-- The logos cant be included into the Package, because their licenses are not clear
%-- as soon as you got the logos name them correctly an uncomment those comments
% \mylogo{ \includegraphics[height=5ex,clip=true , trim=8.1cm  0pt 0pt 0pt ]{logo.eps} }
%\logo{logo.eps}%default Logo at the bottom right

\mylogo{\includegraphics[height=5ex]{/usr/local/texlive/2011/texmf-dist/doc/generic/pstricks/images/tiger.pdf}}
\beamertemplatenavigationsymbolsempty %Navbar on/off




\title[Beamer u. Hohenheim]{Beamer und Hohenheim 00004}
\author[T. G�rlach ]{Tobias G�rlach }

\date[2011 -- 2012]{August 2011 -- }
\institute[Inst.]{Institut}



\begin{document}


 \HohenheimFancyTitle{}



\frame{\tableofcontents}

\frame{\frametitle{Was ist das?} Dieses Dokument zeigt ein Beamertheme mit den Farben der Uni Hohenheim. Ich bezeichne dieses Theme als Hack und m�chte drauf hinweisen, dass er 1. beta ist und 2. nicht unbedingt mit den beamer-Paradigmen �bereinstimmt. Ich �bernehme keinerlei Haftung f�r die vorliegenden Dateien und f�r alles was aus der Benutzung entstehen kann. Zudem m�chte ich darauf hinweisen, dass obschon es sich um die Farben etc., wie in der CI definiert, handelt, dies kein offizielles Theme der Uni ist. Und somit die Uni keinen Support bietet. \bfseries Die Uni hat mit diesem Theme rein gar nichts zu tun. }
\frame{\frametitle{Anleitung zur Anleitung}  
Diese Anleitung ist sowohl Anleitung als auch Beispieldatei. Es wurden keine Logos verwendest, da die Lizenz immer noch nicht klar ist. Diese m�ssen �ber die Uni bezogen werden. Wie diese einzubinden sind, ist dem source-code dieser Datei zu entnehmen.  
 
 }
\frame{\frametitle{Installation} Schiebe die drei Dateien \emph{beamerthemeHohenheim.sty}, \emph{beamercolorthemehohenheim.sty} sowie \emph{beamerouterthemeinfolines\_hohenheim.sty} in einen lokalen Ordner in dem \TeX\ sucht und lasse ggf. texhash durchlaufen. Alternativ kannst du die Dateien auch in den globalen local-texmf Ordner einpflegen. Diesen Ordner kann man mit \texttt{kpsewhich -var-value=TEXMFLOCAL} herausfinden.  Wenn alle D�mme brechen, k�nnen die Dateien auch im Arbeitsverzeichnis abgelegt werden. 

 Vorschl�ge, Bugs, �nderungsw�nsche bitte an \url{tobias.goerlach@uni-hohenheim.de} }
 \frame{\frametitle{Verwendung}
 Zur Verwendung reicht es \texttt{\textbackslash usepackage\{ hobete\} } nach dem Laden von \texttt{beamer} anzugeben. Bei Verwendung von \texttt{beamerposter} sollte die Option \texttt{poster=true} angegeben werden. 
 }
\frame{\frametitle{Farben}
Die Hauptfarbe ist Pantone 2955 C  oder CYMK 100-50-0-45. Das Theme verwendet diese und die mit Wei� geshadete 50\% . Diese Farben sind fest im colortheme verdrahtet und stehen bei Verwendung des Hohenheimer colorthemes zu Verf�gung. \\[2mm]


\begin{tabular}{l>{\ttfamily}ll}
Beschreibung & \textsf{Name} & Test\\
Hauptfarbe&Hohenheim&\fcolorbox{Hohenheim}{Hohenheim}{Test}\\
Hauptfarbe 50\%&Hohenheimbg& \fcolorbox{Hohenheimbg}{Hohenheimbg}{Test}\\
Hauptfarbe f�r BW&HohenheimBW&\fcolorbox{HohenheimBW}{HohenheimBW}{Test}\\
\end{tabular}

}
\frame{\frametitle{Sekund�rfarben}
Diese Farben sollten mit Vorsicht verwendet werden. Vorstellbar ist z.B. eine Verwendung f�r die theorem Umgebeung. Das Theme verwendet diese Farben �berhaupt nicht.\footnote{Einzige Ausnahme: Die Alert-Farbe, z.B. bei Listen.  }
\begin{columns}
\column[T]{.5\textwidth}
Gedeckte Sek Farben\\[1mm]
\begin{tabular}{>{\ttfamily\small}ll}
Hohenheim\_darkgreen &\fcolorbox{Hohenheim_darkgreen}{Hohenheim_darkgreen}{Test}\\
Hohenheim\_green &\fcolorbox{Hohenheim_green}{Hohenheim_green}{Test}\\
Hohenheim\_red &\fcolorbox{Hohenheim_red}{Hohenheim_red}{Test}\\
Hohenheim\_lightblue &\fcolorbox{Hohenheim_lightblue}{Hohenheim_lightblue}{Test}\\
Hohenheim\_midgray &\fcolorbox{Hohenheim_midgray}{Hohenheim_midgray}{Test}\\
\end{tabular}
\column[T]{.5\textwidth}
Leuchtende Sek Farben\\[1mm]
\begin{tabular}{>{\ttfamily\small}ll}
Hohenheim\_glow\_darkgreen &\fcolorbox{Hohenheim_glow_darkgreen}{Hohenheim_glow_darkgreen}{Test}\\
Hohenheim\_glow\_yellow &\fcolorbox{Hohenheim_glow_yellow}{Hohenheim_glow_yellow}{Test}\\
Hohenheim\_glow\_red &\fcolorbox{Hohenheim_glow_red}{Hohenheim_glow_red}{Test}\\
Hohenheim\_glow\_lightblue &\fcolorbox{Hohenheim_glow_lightblue}{Hohenheim_glow_lightblue}{Test}\\
\end{tabular}
\end{columns}

}

\begin{frame}[fragile]
\frametitle{Anwendung der Farben}
Die Farben k�nnen mit den bekannten Befehlen verwendet werden.
\begin{columns}
\column[T]{.5\textwidth} 
\textcolor{Hohenheim_glow_red}{ROT} oder \fcolorbox{Hohenheim_glow_lightblue}{Hohenheim_glow_lightblue}{BOX}%
\column[T]{.5\textwidth}
\begin{lstlisting}[style=tex]
\textcolor{Hohenheim_glow_red}{ROT} oder
 \fcolorbox{Hohenheim_glow_lightblue}
\end{lstlisting}
\end{columns}%
\end{frame}


\begin{frame}[fragile]
 \frametitle{Das Logo} 
Der Befehl \lstinline[style=tex]$ \mylogo{ \includegraphics[height=5ex,clip=true , trim=8.1cm  0pt 0pt 0pt ]{logo.eps} } $  bindet das Logo ein. Dazu muss das Paket graphicx geladen sein da der Befehl dieses verwenden. (Das macht aber beamer schon....) Der \texttt{trim}-Parameter gibt den Zuschnitt des Logos an. Ich empfehle entweder \lstinline[style=tex]$ clip=true , trim=8.1cm  0pt 0pt 0pt$ oder keinen \texttt{trim}. Die H�he sollte nicht ver�ndert werden, da das Layout auf \texttt{5 ex} optimiert ist. Das Logo sollte zu jeder Zeit in Dokument ver�ndert werden k�nnen. 

Dieser Befehl umgeht den \lstinline[style=tex]$\logo$ von beamer. Dieser kann aber weiterhin definiert werden. 

Das Logo erscheint neben dem Titel des Frames. Unbedingt immer den �ber die Uni erh�ltliche .eps File verwenden. Es ist unbedingt das .eps Format zu verwenden, da es sonst zu sehr unsch�nen Render-Schlieren kommt. Bei Verwendung von pdflatex shell-escape anschalten oder mindestens TexLive 2010 verwenden. 
\end{frame}



\begin{frame}[fragile]
\frametitle{Toller Titel...}
... wie geht das? Nun man braucht folgendes:

\begin{lstlisting}[style=tex]
%Header
\usepackage{tikz}
%Document
\frame[plain]{
\begin{tikzpicture}[remember picture,overlay] 
\node [scale=1.7,fill opacity=.2,text opacity=.07] at (current page.south west) {\includegraphics{logoklein.eps} }; 
\node [scale=0.9,fill opacity=.2,text opacity=.8] at (7,0) {\includegraphics{logo.eps} }; 
\end{tikzpicture} 
 \maketitle
 }
\end{lstlisting}
Wobei \texttt{logo.eps} das Logo mit (Wort- u. Bildmarke) und \texttt{logoklein.eps} das Logo ohne den Schriftzug ist. 
\end{frame}
\begin{frame}[fragile]
... oder man definiert einfach 
\begin{lstlisting}[style=tex]
%Header
\HohenheimLogoLang{xy.eps}
\HohenheimLogoKlein{xyz.eps}
%Document
\begin{document}
\HohenheimFancyTitle{<hook>}
\end{lstlisting}
Dabei ist ganz wichtig, dass die Logoangaben VOR \lstinline[style=tex]$\begin{document}$ erfolgen! Es kann, Aufgrund der Arbeitsweise von tikz,  zudem notwendig sein zwei mal zu Kompilieren. \texttt{<hook>} kann mit beliebigen weiterem Inhalt gef�llt werden, zB. ein Institutslogo. 
\end{frame}

\begin{frame}[fragile]
\frametitle{Sectionpage}
Die Sectionpage, \lstinline[style=tex]$\sectionpage$ , gibt die aktuelle Section auf einer separaten Seite aus. Die Optionen sind:\\
\lstinline[style=tex]$\sectionpage[plain=true]$ : Sectionpage ohne Kopf- und Fusszeile.  
\end{frame}



\begin{frame} %
\frametitle{Framelist}

\printframelist
% Copyright 2010-2012  by tgoerlach < tobias.goerlach@uni-hohenheim.de >
% This is a example file for beamerthemehohenheim
% https://bitbucket.org/tobig/hohenheimbeamertheme/
%
% This file may be distributed and/or modified
%
% 1. under the LaTeX Project Public License and/or
% 2. under the GNU Public License.
%
% See the file doc/licenses/LICENSE for more details.


\documentclass[xcolor={table}]{beamer}

\usepackage[ngerman]{babel}
\usepackage[latin1]{inputenc}
\usepackage[T1]{fontenc}
\usepackage{lmodern}


% Document specific stuff ()
\usepackage{etex}% in some cases that helps% eg if you get those no room for newdimen errors
% here we need it for pgfplots
\usepackage{booktabs, url, calc}
\usepackage{phaistos}% Just for the bee
\newcommand{\bee}{\scalebox{0.5}{\PHbee}\ }
\usepackage{tabularx} % tabu is even better
\usepackage{amsmath, amssymb} % ok no commend needed here
\usepackage{listings}% for code % note that beamer wants you to put code only into fragile frames (ref to the beamerusersguide)
\lstdefinestyle{tex}{language=tex, basicstyle={\footnotesize\ttfamily} , breaklines=true}
\usepackage{pgfplots}% best function plotter for tex
 \pgfplotsset{compat=newest}
\pgfplotsset{tick label style={font=\small\sffamily}, label style={font=\small\sffamily}, legend style={font=\footnotesize\sffamily},
}







%-- Theme Stuff
\usepackage[framelist front={o\ }, framelist back={\ }]{hobete}
%--that all what you have to do. 


%\HohenheimLogoKlein{/usr/local/texlive/2011/texmf-dist/doc/generic/pstricks/images/tiger.pdf} %Use the eps provided by the University
%\HohenheimLogoLang{/usr/local/texlive/2011/texmf-dist/doc/generic/pstricks/images/tiger.pdf} % Lang indicates that here mus be insert the "Wort und Bildmarke"
%Klein indicates, that here should only the Bildmarke be inserted. 
%-- You can gain them from the University CI Center
%-- The logos cant be included into the Package, because their licenses are not clear
%-- as soon as you got the logos name them correctly an uncomment those comments
% \mylogo{ \includegraphics[height=5ex,clip=true , trim=8.1cm  0pt 0pt 0pt ]{logo.eps} }
%\logo{logo.eps}%default Logo at the bottom right

\mylogo{\includegraphics[height=5ex]{/usr/local/texlive/2011/texmf-dist/doc/generic/pstricks/images/tiger.pdf}}
\beamertemplatenavigationsymbolsempty %Navbar on/off




\title[Beamer u. Hohenheim]{Beamer und Hohenheim 00004}
\author[T. G�rlach ]{Tobias G�rlach }

\date[2011 -- 2012]{August 2011 -- }
\institute[Inst.]{Institut}



\begin{document}


 \HohenheimFancyTitle{}



\frame{\tableofcontents}

\frame{\frametitle{Was ist das?} Dieses Dokument zeigt ein Beamertheme mit den Farben der Uni Hohenheim. Ich bezeichne dieses Theme als Hack und m�chte drauf hinweisen, dass er 1. beta ist und 2. nicht unbedingt mit den beamer-Paradigmen �bereinstimmt. Ich �bernehme keinerlei Haftung f�r die vorliegenden Dateien und f�r alles was aus der Benutzung entstehen kann. Zudem m�chte ich darauf hinweisen, dass obschon es sich um die Farben etc., wie in der CI definiert, handelt, dies kein offizielles Theme der Uni ist. Und somit die Uni keinen Support bietet. \bfseries Die Uni hat mit diesem Theme rein gar nichts zu tun. }
\frame{\frametitle{Anleitung zur Anleitung}  
Diese Anleitung ist sowohl Anleitung als auch Beispieldatei. Es wurden keine Logos verwendest, da die Lizenz immer noch nicht klar ist. Diese m�ssen �ber die Uni bezogen werden. Wie diese einzubinden sind, ist dem source-code dieser Datei zu entnehmen.  
 
 }
\frame{\frametitle{Installation} Schiebe die drei Dateien \emph{beamerthemeHohenheim.sty}, \emph{beamercolorthemehohenheim.sty} sowie \emph{beamerouterthemeinfolines\_hohenheim.sty} in einen lokalen Ordner in dem \TeX\ sucht und lasse ggf. texhash durchlaufen. Alternativ kannst du die Dateien auch in den globalen local-texmf Ordner einpflegen. Diesen Ordner kann man mit \texttt{kpsewhich -var-value=TEXMFLOCAL} herausfinden.  Wenn alle D�mme brechen, k�nnen die Dateien auch im Arbeitsverzeichnis abgelegt werden. 

 Vorschl�ge, Bugs, �nderungsw�nsche bitte an \url{tobias.goerlach@uni-hohenheim.de} }
 \frame{\frametitle{Verwendung}
 Zur Verwendung reicht es \texttt{\textbackslash usepackage\{ hobete\} } nach dem Laden von \texttt{beamer} anzugeben. Bei Verwendung von \texttt{beamerposter} sollte die Option \texttt{poster=true} angegeben werden. 
 }
\frame{\frametitle{Farben}
Die Hauptfarbe ist Pantone 2955 C  oder CYMK 100-50-0-45. Das Theme verwendet diese und die mit Wei� geshadete 50\% . Diese Farben sind fest im colortheme verdrahtet und stehen bei Verwendung des Hohenheimer colorthemes zu Verf�gung. \\[2mm]


\begin{tabular}{l>{\ttfamily}ll}
Beschreibung & \textsf{Name} & Test\\
Hauptfarbe&Hohenheim&\fcolorbox{Hohenheim}{Hohenheim}{Test}\\
Hauptfarbe 50\%&Hohenheimbg& \fcolorbox{Hohenheimbg}{Hohenheimbg}{Test}\\
Hauptfarbe f�r BW&HohenheimBW&\fcolorbox{HohenheimBW}{HohenheimBW}{Test}\\
\end{tabular}

}
\frame{\frametitle{Sekund�rfarben}
Diese Farben sollten mit Vorsicht verwendet werden. Vorstellbar ist z.B. eine Verwendung f�r die theorem Umgebeung. Das Theme verwendet diese Farben �berhaupt nicht.\footnote{Einzige Ausnahme: Die Alert-Farbe, z.B. bei Listen.  }
\begin{columns}
\column[T]{.5\textwidth}
Gedeckte Sek Farben\\[1mm]
\begin{tabular}{>{\ttfamily\small}ll}
Hohenheim\_darkgreen &\fcolorbox{Hohenheim_darkgreen}{Hohenheim_darkgreen}{Test}\\
Hohenheim\_green &\fcolorbox{Hohenheim_green}{Hohenheim_green}{Test}\\
Hohenheim\_red &\fcolorbox{Hohenheim_red}{Hohenheim_red}{Test}\\
Hohenheim\_lightblue &\fcolorbox{Hohenheim_lightblue}{Hohenheim_lightblue}{Test}\\
Hohenheim\_midgray &\fcolorbox{Hohenheim_midgray}{Hohenheim_midgray}{Test}\\
\end{tabular}
\column[T]{.5\textwidth}
Leuchtende Sek Farben\\[1mm]
\begin{tabular}{>{\ttfamily\small}ll}
Hohenheim\_glow\_darkgreen &\fcolorbox{Hohenheim_glow_darkgreen}{Hohenheim_glow_darkgreen}{Test}\\
Hohenheim\_glow\_yellow &\fcolorbox{Hohenheim_glow_yellow}{Hohenheim_glow_yellow}{Test}\\
Hohenheim\_glow\_red &\fcolorbox{Hohenheim_glow_red}{Hohenheim_glow_red}{Test}\\
Hohenheim\_glow\_lightblue &\fcolorbox{Hohenheim_glow_lightblue}{Hohenheim_glow_lightblue}{Test}\\
\end{tabular}
\end{columns}

}

\begin{frame}[fragile]
\frametitle{Anwendung der Farben}
Die Farben k�nnen mit den bekannten Befehlen verwendet werden.
\begin{columns}
\column[T]{.5\textwidth} 
\textcolor{Hohenheim_glow_red}{ROT} oder \fcolorbox{Hohenheim_glow_lightblue}{Hohenheim_glow_lightblue}{BOX}%
\column[T]{.5\textwidth}
\begin{lstlisting}[style=tex]
\textcolor{Hohenheim_glow_red}{ROT} oder
 \fcolorbox{Hohenheim_glow_lightblue}
\end{lstlisting}
\end{columns}%
\end{frame}


\begin{frame}[fragile]
 \frametitle{Das Logo} 
Der Befehl \lstinline[style=tex]$ \mylogo{ \includegraphics[height=5ex,clip=true , trim=8.1cm  0pt 0pt 0pt ]{logo.eps} } $  bindet das Logo ein. Dazu muss das Paket graphicx geladen sein da der Befehl dieses verwenden. (Das macht aber beamer schon....) Der \texttt{trim}-Parameter gibt den Zuschnitt des Logos an. Ich empfehle entweder \lstinline[style=tex]$ clip=true , trim=8.1cm  0pt 0pt 0pt$ oder keinen \texttt{trim}. Die H�he sollte nicht ver�ndert werden, da das Layout auf \texttt{5 ex} optimiert ist. Das Logo sollte zu jeder Zeit in Dokument ver�ndert werden k�nnen. 

Dieser Befehl umgeht den \lstinline[style=tex]$\logo$ von beamer. Dieser kann aber weiterhin definiert werden. 

Das Logo erscheint neben dem Titel des Frames. Unbedingt immer den �ber die Uni erh�ltliche .eps File verwenden. Es ist unbedingt das .eps Format zu verwenden, da es sonst zu sehr unsch�nen Render-Schlieren kommt. Bei Verwendung von pdflatex shell-escape anschalten oder mindestens TexLive 2010 verwenden. 
\end{frame}



\begin{frame}[fragile]
\frametitle{Toller Titel...}
... wie geht das? Nun man braucht folgendes:

\begin{lstlisting}[style=tex]
%Header
\usepackage{tikz}
%Document
\frame[plain]{
\begin{tikzpicture}[remember picture,overlay] 
\node [scale=1.7,fill opacity=.2,text opacity=.07] at (current page.south west) {\includegraphics{logoklein.eps} }; 
\node [scale=0.9,fill opacity=.2,text opacity=.8] at (7,0) {\includegraphics{logo.eps} }; 
\end{tikzpicture} 
 \maketitle
 }
\end{lstlisting}
Wobei \texttt{logo.eps} das Logo mit (Wort- u. Bildmarke) und \texttt{logoklein.eps} das Logo ohne den Schriftzug ist. 
\end{frame}
\begin{frame}[fragile]
... oder man definiert einfach 
\begin{lstlisting}[style=tex]
%Header
\HohenheimLogoLang{xy.eps}
\HohenheimLogoKlein{xyz.eps}
%Document
\begin{document}
\HohenheimFancyTitle{<hook>}
\end{lstlisting}
Dabei ist ganz wichtig, dass die Logoangaben VOR \lstinline[style=tex]$\begin{document}$ erfolgen! Es kann, Aufgrund der Arbeitsweise von tikz,  zudem notwendig sein zwei mal zu Kompilieren. \texttt{<hook>} kann mit beliebigen weiterem Inhalt gef�llt werden, zB. ein Institutslogo. 
\end{frame}

\begin{frame}[fragile]
\frametitle{Sectionpage}
Die Sectionpage, \lstinline[style=tex]$\sectionpage$ , gibt die aktuelle Section auf einer separaten Seite aus. Die Optionen sind:\\
\lstinline[style=tex]$\sectionpage[plain=true]$ : Sectionpage ohne Kopf- und Fusszeile.  
\end{frame}



\begin{frame} %
\frametitle{Framelist}

\printframelist
% Copyright 2010-2012  by tgoerlach < tobias.goerlach@uni-hohenheim.de >
% This is a example file for beamerthemehohenheim
% https://bitbucket.org/tobig/hohenheimbeamertheme/
%
% This file may be distributed and/or modified
%
% 1. under the LaTeX Project Public License and/or
% 2. under the GNU Public License.
%
% See the file doc/licenses/LICENSE for more details.


\documentclass[xcolor={table}]{beamer}

\usepackage[ngerman]{babel}
\usepackage[latin1]{inputenc}
\usepackage[T1]{fontenc}
\usepackage{lmodern}


% Document specific stuff ()
\usepackage{etex}% in some cases that helps% eg if you get those no room for newdimen errors
% here we need it for pgfplots
\usepackage{booktabs, url, calc}
\usepackage{phaistos}% Just for the bee
\newcommand{\bee}{\scalebox{0.5}{\PHbee}\ }
\usepackage{tabularx} % tabu is even better
\usepackage{amsmath, amssymb} % ok no commend needed here
\usepackage{listings}% for code % note that beamer wants you to put code only into fragile frames (ref to the beamerusersguide)
\lstdefinestyle{tex}{language=tex, basicstyle={\footnotesize\ttfamily} , breaklines=true}
\usepackage{pgfplots}% best function plotter for tex
 \pgfplotsset{compat=newest}
\pgfplotsset{tick label style={font=\small\sffamily}, label style={font=\small\sffamily}, legend style={font=\footnotesize\sffamily},
}







%-- Theme Stuff
\usepackage[framelist front={o\ }, framelist back={\ }]{hobete}
%--that all what you have to do. 


%\HohenheimLogoKlein{/usr/local/texlive/2011/texmf-dist/doc/generic/pstricks/images/tiger.pdf} %Use the eps provided by the University
%\HohenheimLogoLang{/usr/local/texlive/2011/texmf-dist/doc/generic/pstricks/images/tiger.pdf} % Lang indicates that here mus be insert the "Wort und Bildmarke"
%Klein indicates, that here should only the Bildmarke be inserted. 
%-- You can gain them from the University CI Center
%-- The logos cant be included into the Package, because their licenses are not clear
%-- as soon as you got the logos name them correctly an uncomment those comments
% \mylogo{ \includegraphics[height=5ex,clip=true , trim=8.1cm  0pt 0pt 0pt ]{logo.eps} }
%\logo{logo.eps}%default Logo at the bottom right

\mylogo{\includegraphics[height=5ex]{/usr/local/texlive/2011/texmf-dist/doc/generic/pstricks/images/tiger.pdf}}
\beamertemplatenavigationsymbolsempty %Navbar on/off




\title[Beamer u. Hohenheim]{Beamer und Hohenheim 00004}
\author[T. G�rlach ]{Tobias G�rlach }

\date[2011 -- 2012]{August 2011 -- }
\institute[Inst.]{Institut}



\begin{document}


 \HohenheimFancyTitle{}



\frame{\tableofcontents}

\frame{\frametitle{Was ist das?} Dieses Dokument zeigt ein Beamertheme mit den Farben der Uni Hohenheim. Ich bezeichne dieses Theme als Hack und m�chte drauf hinweisen, dass er 1. beta ist und 2. nicht unbedingt mit den beamer-Paradigmen �bereinstimmt. Ich �bernehme keinerlei Haftung f�r die vorliegenden Dateien und f�r alles was aus der Benutzung entstehen kann. Zudem m�chte ich darauf hinweisen, dass obschon es sich um die Farben etc., wie in der CI definiert, handelt, dies kein offizielles Theme der Uni ist. Und somit die Uni keinen Support bietet. \bfseries Die Uni hat mit diesem Theme rein gar nichts zu tun. }
\frame{\frametitle{Anleitung zur Anleitung}  
Diese Anleitung ist sowohl Anleitung als auch Beispieldatei. Es wurden keine Logos verwendest, da die Lizenz immer noch nicht klar ist. Diese m�ssen �ber die Uni bezogen werden. Wie diese einzubinden sind, ist dem source-code dieser Datei zu entnehmen.  
 
 }
\frame{\frametitle{Installation} Schiebe die drei Dateien \emph{beamerthemeHohenheim.sty}, \emph{beamercolorthemehohenheim.sty} sowie \emph{beamerouterthemeinfolines\_hohenheim.sty} in einen lokalen Ordner in dem \TeX\ sucht und lasse ggf. texhash durchlaufen. Alternativ kannst du die Dateien auch in den globalen local-texmf Ordner einpflegen. Diesen Ordner kann man mit \texttt{kpsewhich -var-value=TEXMFLOCAL} herausfinden.  Wenn alle D�mme brechen, k�nnen die Dateien auch im Arbeitsverzeichnis abgelegt werden. 

 Vorschl�ge, Bugs, �nderungsw�nsche bitte an \url{tobias.goerlach@uni-hohenheim.de} }
 \frame{\frametitle{Verwendung}
 Zur Verwendung reicht es \texttt{\textbackslash usepackage\{ hobete\} } nach dem Laden von \texttt{beamer} anzugeben. Bei Verwendung von \texttt{beamerposter} sollte die Option \texttt{poster=true} angegeben werden. 
 }
\frame{\frametitle{Farben}
Die Hauptfarbe ist Pantone 2955 C  oder CYMK 100-50-0-45. Das Theme verwendet diese und die mit Wei� geshadete 50\% . Diese Farben sind fest im colortheme verdrahtet und stehen bei Verwendung des Hohenheimer colorthemes zu Verf�gung. \\[2mm]


\begin{tabular}{l>{\ttfamily}ll}
Beschreibung & \textsf{Name} & Test\\
Hauptfarbe&Hohenheim&\fcolorbox{Hohenheim}{Hohenheim}{Test}\\
Hauptfarbe 50\%&Hohenheimbg& \fcolorbox{Hohenheimbg}{Hohenheimbg}{Test}\\
Hauptfarbe f�r BW&HohenheimBW&\fcolorbox{HohenheimBW}{HohenheimBW}{Test}\\
\end{tabular}

}
\frame{\frametitle{Sekund�rfarben}
Diese Farben sollten mit Vorsicht verwendet werden. Vorstellbar ist z.B. eine Verwendung f�r die theorem Umgebeung. Das Theme verwendet diese Farben �berhaupt nicht.\footnote{Einzige Ausnahme: Die Alert-Farbe, z.B. bei Listen.  }
\begin{columns}
\column[T]{.5\textwidth}
Gedeckte Sek Farben\\[1mm]
\begin{tabular}{>{\ttfamily\small}ll}
Hohenheim\_darkgreen &\fcolorbox{Hohenheim_darkgreen}{Hohenheim_darkgreen}{Test}\\
Hohenheim\_green &\fcolorbox{Hohenheim_green}{Hohenheim_green}{Test}\\
Hohenheim\_red &\fcolorbox{Hohenheim_red}{Hohenheim_red}{Test}\\
Hohenheim\_lightblue &\fcolorbox{Hohenheim_lightblue}{Hohenheim_lightblue}{Test}\\
Hohenheim\_midgray &\fcolorbox{Hohenheim_midgray}{Hohenheim_midgray}{Test}\\
\end{tabular}
\column[T]{.5\textwidth}
Leuchtende Sek Farben\\[1mm]
\begin{tabular}{>{\ttfamily\small}ll}
Hohenheim\_glow\_darkgreen &\fcolorbox{Hohenheim_glow_darkgreen}{Hohenheim_glow_darkgreen}{Test}\\
Hohenheim\_glow\_yellow &\fcolorbox{Hohenheim_glow_yellow}{Hohenheim_glow_yellow}{Test}\\
Hohenheim\_glow\_red &\fcolorbox{Hohenheim_glow_red}{Hohenheim_glow_red}{Test}\\
Hohenheim\_glow\_lightblue &\fcolorbox{Hohenheim_glow_lightblue}{Hohenheim_glow_lightblue}{Test}\\
\end{tabular}
\end{columns}

}

\begin{frame}[fragile]
\frametitle{Anwendung der Farben}
Die Farben k�nnen mit den bekannten Befehlen verwendet werden.
\begin{columns}
\column[T]{.5\textwidth} 
\textcolor{Hohenheim_glow_red}{ROT} oder \fcolorbox{Hohenheim_glow_lightblue}{Hohenheim_glow_lightblue}{BOX}%
\column[T]{.5\textwidth}
\begin{lstlisting}[style=tex]
\textcolor{Hohenheim_glow_red}{ROT} oder
 \fcolorbox{Hohenheim_glow_lightblue}
\end{lstlisting}
\end{columns}%
\end{frame}


\begin{frame}[fragile]
 \frametitle{Das Logo} 
Der Befehl \lstinline[style=tex]$ \mylogo{ \includegraphics[height=5ex,clip=true , trim=8.1cm  0pt 0pt 0pt ]{logo.eps} } $  bindet das Logo ein. Dazu muss das Paket graphicx geladen sein da der Befehl dieses verwenden. (Das macht aber beamer schon....) Der \texttt{trim}-Parameter gibt den Zuschnitt des Logos an. Ich empfehle entweder \lstinline[style=tex]$ clip=true , trim=8.1cm  0pt 0pt 0pt$ oder keinen \texttt{trim}. Die H�he sollte nicht ver�ndert werden, da das Layout auf \texttt{5 ex} optimiert ist. Das Logo sollte zu jeder Zeit in Dokument ver�ndert werden k�nnen. 

Dieser Befehl umgeht den \lstinline[style=tex]$\logo$ von beamer. Dieser kann aber weiterhin definiert werden. 

Das Logo erscheint neben dem Titel des Frames. Unbedingt immer den �ber die Uni erh�ltliche .eps File verwenden. Es ist unbedingt das .eps Format zu verwenden, da es sonst zu sehr unsch�nen Render-Schlieren kommt. Bei Verwendung von pdflatex shell-escape anschalten oder mindestens TexLive 2010 verwenden. 
\end{frame}



\begin{frame}[fragile]
\frametitle{Toller Titel...}
... wie geht das? Nun man braucht folgendes:

\begin{lstlisting}[style=tex]
%Header
\usepackage{tikz}
%Document
\frame[plain]{
\begin{tikzpicture}[remember picture,overlay] 
\node [scale=1.7,fill opacity=.2,text opacity=.07] at (current page.south west) {\includegraphics{logoklein.eps} }; 
\node [scale=0.9,fill opacity=.2,text opacity=.8] at (7,0) {\includegraphics{logo.eps} }; 
\end{tikzpicture} 
 \maketitle
 }
\end{lstlisting}
Wobei \texttt{logo.eps} das Logo mit (Wort- u. Bildmarke) und \texttt{logoklein.eps} das Logo ohne den Schriftzug ist. 
\end{frame}
\begin{frame}[fragile]
... oder man definiert einfach 
\begin{lstlisting}[style=tex]
%Header
\HohenheimLogoLang{xy.eps}
\HohenheimLogoKlein{xyz.eps}
%Document
\begin{document}
\HohenheimFancyTitle{<hook>}
\end{lstlisting}
Dabei ist ganz wichtig, dass die Logoangaben VOR \lstinline[style=tex]$\begin{document}$ erfolgen! Es kann, Aufgrund der Arbeitsweise von tikz,  zudem notwendig sein zwei mal zu Kompilieren. \texttt{<hook>} kann mit beliebigen weiterem Inhalt gef�llt werden, zB. ein Institutslogo. 
\end{frame}

\begin{frame}[fragile]
\frametitle{Sectionpage}
Die Sectionpage, \lstinline[style=tex]$\sectionpage$ , gibt die aktuelle Section auf einer separaten Seite aus. Die Optionen sind:\\
\lstinline[style=tex]$\sectionpage[plain=true]$ : Sectionpage ohne Kopf- und Fusszeile.  
\end{frame}



\begin{frame} %
\frametitle{Framelist}

\printframelist
% Copyright 2010-2012  by tgoerlach < tobias.goerlach@uni-hohenheim.de >
% This is a example file for beamerthemehohenheim
% https://bitbucket.org/tobig/hohenheimbeamertheme/
%
% This file may be distributed and/or modified
%
% 1. under the LaTeX Project Public License and/or
% 2. under the GNU Public License.
%
% See the file doc/licenses/LICENSE for more details.


\documentclass[xcolor={table}]{beamer}

\usepackage[ngerman]{babel}
\usepackage[latin1]{inputenc}
\usepackage[T1]{fontenc}
\usepackage{lmodern}


% Document specific stuff ()
\usepackage{etex}% in some cases that helps% eg if you get those no room for newdimen errors
% here we need it for pgfplots
\usepackage{booktabs, url, calc}
\usepackage{phaistos}% Just for the bee
\newcommand{\bee}{\scalebox{0.5}{\PHbee}\ }
\usepackage{tabularx} % tabu is even better
\usepackage{amsmath, amssymb} % ok no commend needed here
\usepackage{listings}% for code % note that beamer wants you to put code only into fragile frames (ref to the beamerusersguide)
\lstdefinestyle{tex}{language=tex, basicstyle={\footnotesize\ttfamily} , breaklines=true}
\usepackage{pgfplots}% best function plotter for tex
 \pgfplotsset{compat=newest}
\pgfplotsset{tick label style={font=\small\sffamily}, label style={font=\small\sffamily}, legend style={font=\footnotesize\sffamily},
}







%-- Theme Stuff
\usepackage[framelist front={o\ }, framelist back={\ }]{hobete}
%--that all what you have to do. 


%\HohenheimLogoKlein{/usr/local/texlive/2011/texmf-dist/doc/generic/pstricks/images/tiger.pdf} %Use the eps provided by the University
%\HohenheimLogoLang{/usr/local/texlive/2011/texmf-dist/doc/generic/pstricks/images/tiger.pdf} % Lang indicates that here mus be insert the "Wort und Bildmarke"
%Klein indicates, that here should only the Bildmarke be inserted. 
%-- You can gain them from the University CI Center
%-- The logos cant be included into the Package, because their licenses are not clear
%-- as soon as you got the logos name them correctly an uncomment those comments
% \mylogo{ \includegraphics[height=5ex,clip=true , trim=8.1cm  0pt 0pt 0pt ]{logo.eps} }
%\logo{logo.eps}%default Logo at the bottom right

\mylogo{\includegraphics[height=5ex]{/usr/local/texlive/2011/texmf-dist/doc/generic/pstricks/images/tiger.pdf}}
\beamertemplatenavigationsymbolsempty %Navbar on/off




\title[Beamer u. Hohenheim]{Beamer und Hohenheim 00004}
\author[T. G�rlach ]{Tobias G�rlach }

\date[2011 -- 2012]{August 2011 -- }
\institute[Inst.]{Institut}



\begin{document}


 \HohenheimFancyTitle{}



\frame{\tableofcontents}

\frame{\frametitle{Was ist das?} Dieses Dokument zeigt ein Beamertheme mit den Farben der Uni Hohenheim. Ich bezeichne dieses Theme als Hack und m�chte drauf hinweisen, dass er 1. beta ist und 2. nicht unbedingt mit den beamer-Paradigmen �bereinstimmt. Ich �bernehme keinerlei Haftung f�r die vorliegenden Dateien und f�r alles was aus der Benutzung entstehen kann. Zudem m�chte ich darauf hinweisen, dass obschon es sich um die Farben etc., wie in der CI definiert, handelt, dies kein offizielles Theme der Uni ist. Und somit die Uni keinen Support bietet. \bfseries Die Uni hat mit diesem Theme rein gar nichts zu tun. }
\frame{\frametitle{Anleitung zur Anleitung}  
Diese Anleitung ist sowohl Anleitung als auch Beispieldatei. Es wurden keine Logos verwendest, da die Lizenz immer noch nicht klar ist. Diese m�ssen �ber die Uni bezogen werden. Wie diese einzubinden sind, ist dem source-code dieser Datei zu entnehmen.  
 
 }
\frame{\frametitle{Installation} Schiebe die drei Dateien \emph{beamerthemeHohenheim.sty}, \emph{beamercolorthemehohenheim.sty} sowie \emph{beamerouterthemeinfolines\_hohenheim.sty} in einen lokalen Ordner in dem \TeX\ sucht und lasse ggf. texhash durchlaufen. Alternativ kannst du die Dateien auch in den globalen local-texmf Ordner einpflegen. Diesen Ordner kann man mit \texttt{kpsewhich -var-value=TEXMFLOCAL} herausfinden.  Wenn alle D�mme brechen, k�nnen die Dateien auch im Arbeitsverzeichnis abgelegt werden. 

 Vorschl�ge, Bugs, �nderungsw�nsche bitte an \url{tobias.goerlach@uni-hohenheim.de} }
 \frame{\frametitle{Verwendung}
 Zur Verwendung reicht es \texttt{\textbackslash usepackage\{ hobete\} } nach dem Laden von \texttt{beamer} anzugeben. Bei Verwendung von \texttt{beamerposter} sollte die Option \texttt{poster=true} angegeben werden. 
 }
\frame{\frametitle{Farben}
Die Hauptfarbe ist Pantone 2955 C  oder CYMK 100-50-0-45. Das Theme verwendet diese und die mit Wei� geshadete 50\% . Diese Farben sind fest im colortheme verdrahtet und stehen bei Verwendung des Hohenheimer colorthemes zu Verf�gung. \\[2mm]


\begin{tabular}{l>{\ttfamily}ll}
Beschreibung & \textsf{Name} & Test\\
Hauptfarbe&Hohenheim&\fcolorbox{Hohenheim}{Hohenheim}{Test}\\
Hauptfarbe 50\%&Hohenheimbg& \fcolorbox{Hohenheimbg}{Hohenheimbg}{Test}\\
Hauptfarbe f�r BW&HohenheimBW&\fcolorbox{HohenheimBW}{HohenheimBW}{Test}\\
\end{tabular}

}
\frame{\frametitle{Sekund�rfarben}
Diese Farben sollten mit Vorsicht verwendet werden. Vorstellbar ist z.B. eine Verwendung f�r die theorem Umgebeung. Das Theme verwendet diese Farben �berhaupt nicht.\footnote{Einzige Ausnahme: Die Alert-Farbe, z.B. bei Listen.  }
\begin{columns}
\column[T]{.5\textwidth}
Gedeckte Sek Farben\\[1mm]
\begin{tabular}{>{\ttfamily\small}ll}
Hohenheim\_darkgreen &\fcolorbox{Hohenheim_darkgreen}{Hohenheim_darkgreen}{Test}\\
Hohenheim\_green &\fcolorbox{Hohenheim_green}{Hohenheim_green}{Test}\\
Hohenheim\_red &\fcolorbox{Hohenheim_red}{Hohenheim_red}{Test}\\
Hohenheim\_lightblue &\fcolorbox{Hohenheim_lightblue}{Hohenheim_lightblue}{Test}\\
Hohenheim\_midgray &\fcolorbox{Hohenheim_midgray}{Hohenheim_midgray}{Test}\\
\end{tabular}
\column[T]{.5\textwidth}
Leuchtende Sek Farben\\[1mm]
\begin{tabular}{>{\ttfamily\small}ll}
Hohenheim\_glow\_darkgreen &\fcolorbox{Hohenheim_glow_darkgreen}{Hohenheim_glow_darkgreen}{Test}\\
Hohenheim\_glow\_yellow &\fcolorbox{Hohenheim_glow_yellow}{Hohenheim_glow_yellow}{Test}\\
Hohenheim\_glow\_red &\fcolorbox{Hohenheim_glow_red}{Hohenheim_glow_red}{Test}\\
Hohenheim\_glow\_lightblue &\fcolorbox{Hohenheim_glow_lightblue}{Hohenheim_glow_lightblue}{Test}\\
\end{tabular}
\end{columns}

}

\begin{frame}[fragile]
\frametitle{Anwendung der Farben}
Die Farben k�nnen mit den bekannten Befehlen verwendet werden.
\begin{columns}
\column[T]{.5\textwidth} 
\textcolor{Hohenheim_glow_red}{ROT} oder \fcolorbox{Hohenheim_glow_lightblue}{Hohenheim_glow_lightblue}{BOX}%
\column[T]{.5\textwidth}
\begin{lstlisting}[style=tex]
\textcolor{Hohenheim_glow_red}{ROT} oder
 \fcolorbox{Hohenheim_glow_lightblue}
\end{lstlisting}
\end{columns}%
\end{frame}


\begin{frame}[fragile]
 \frametitle{Das Logo} 
Der Befehl \lstinline[style=tex]$ \mylogo{ \includegraphics[height=5ex,clip=true , trim=8.1cm  0pt 0pt 0pt ]{logo.eps} } $  bindet das Logo ein. Dazu muss das Paket graphicx geladen sein da der Befehl dieses verwenden. (Das macht aber beamer schon....) Der \texttt{trim}-Parameter gibt den Zuschnitt des Logos an. Ich empfehle entweder \lstinline[style=tex]$ clip=true , trim=8.1cm  0pt 0pt 0pt$ oder keinen \texttt{trim}. Die H�he sollte nicht ver�ndert werden, da das Layout auf \texttt{5 ex} optimiert ist. Das Logo sollte zu jeder Zeit in Dokument ver�ndert werden k�nnen. 

Dieser Befehl umgeht den \lstinline[style=tex]$\logo$ von beamer. Dieser kann aber weiterhin definiert werden. 

Das Logo erscheint neben dem Titel des Frames. Unbedingt immer den �ber die Uni erh�ltliche .eps File verwenden. Es ist unbedingt das .eps Format zu verwenden, da es sonst zu sehr unsch�nen Render-Schlieren kommt. Bei Verwendung von pdflatex shell-escape anschalten oder mindestens TexLive 2010 verwenden. 
\end{frame}



\begin{frame}[fragile]
\frametitle{Toller Titel...}
... wie geht das? Nun man braucht folgendes:

\begin{lstlisting}[style=tex]
%Header
\usepackage{tikz}
%Document
\frame[plain]{
\begin{tikzpicture}[remember picture,overlay] 
\node [scale=1.7,fill opacity=.2,text opacity=.07] at (current page.south west) {\includegraphics{logoklein.eps} }; 
\node [scale=0.9,fill opacity=.2,text opacity=.8] at (7,0) {\includegraphics{logo.eps} }; 
\end{tikzpicture} 
 \maketitle
 }
\end{lstlisting}
Wobei \texttt{logo.eps} das Logo mit (Wort- u. Bildmarke) und \texttt{logoklein.eps} das Logo ohne den Schriftzug ist. 
\end{frame}
\begin{frame}[fragile]
... oder man definiert einfach 
\begin{lstlisting}[style=tex]
%Header
\HohenheimLogoLang{xy.eps}
\HohenheimLogoKlein{xyz.eps}
%Document
\begin{document}
\HohenheimFancyTitle{<hook>}
\end{lstlisting}
Dabei ist ganz wichtig, dass die Logoangaben VOR \lstinline[style=tex]$\begin{document}$ erfolgen! Es kann, Aufgrund der Arbeitsweise von tikz,  zudem notwendig sein zwei mal zu Kompilieren. \texttt{<hook>} kann mit beliebigen weiterem Inhalt gef�llt werden, zB. ein Institutslogo. 
\end{frame}

\begin{frame}[fragile]
\frametitle{Sectionpage}
Die Sectionpage, \lstinline[style=tex]$\sectionpage$ , gibt die aktuelle Section auf einer separaten Seite aus. Die Optionen sind:\\
\lstinline[style=tex]$\sectionpage[plain=true]$ : Sectionpage ohne Kopf- und Fusszeile.  
\end{frame}



\begin{frame} %
\frametitle{Framelist}

\printframelist
\input{./hobete_doc.frml}
\end{frame}

\begin{frame}[fragile]\frametitle{Die Framelist}
Die Framelist ist eine Liste der Frametitles. Diese wird �hnlich zu toc�s in einer Hilfsdatei mit der Endung \texttt{.frml} abgelegt. Der Befehl \lstinline[style=tex]$ \printframelist $ gibt die Liste aus. Mit der Klassenoption \lstinline[style=tex]$ framelist front$ und \lstinline[style=tex]$ framelist back $ k�nnen die Token eingestellt werden, die vor bzw nach den Eintr�gen erscheinen. Defaultwerte sind \lstinline[style=tex]{ $\cdot$ } und \lstinline[style=tex]$ \\ $ .
\end{frame}


\newcommand{\p}{{\footnotesize +}\ }\newcommand{\pp}{{\footnotesize ++}\ }

\frame{\frametitle{Change-Log}\scriptsize
\begin{tabularx}{\textwidth}{XXX}
Version & Changes & Status-Name\\
0001& \p Colortheme\newline \p Fancy Framtitle\newline \p mylogo-Command  &  beta-hack\\[15mm]
0001a & \bee   improved frametitle\newline \bee correct name for yellow \newline \p better  Logo control \newline \p added \texttt{HohenheimFancyTitle} &\\[15mm]
0001b & \bee   minor bug fixes\newline &\\
0001c & \p sectionpage now with plain option \p Better \texttt{HohenheimFancyTitle} \p altered item labels & \\
0002a & \p added framelist   & beta-hack release candiate (the steped version number indicates that this is a realease version) \\
\end{tabularx}\vfill
\tiny
\p = new feature, \pp = significantly improved feature,  - = removed feature, \bee  bug-fix
}
\frame{\frametitle{Change-Log Vorts.}\scriptsize
\begin{tabularx}{\textwidth}{XXX}
Version & Changes & Status-Name\\
0002b & \p added Support for beamerposter & \\
0003 & \pp tidy up of the code, from now on its a real Package & \\
0003a & \pp growing support for beamerposter & \\
0004 & \pp latex3 plus options & \\
\end{tabularx}\vfill
\tiny
\p = new feature, \pp = significantly improved feature,  - = removed feature, \bee  bug-fix
}





\frame{\frametitle{To Do }\scriptsize

\begin{tabularx}{\textwidth}{XXX}
Importance & Changes & Progress/Probs\\
0& \p Support for beamerarticle &  non \\ 
\rowcolor{black!10} 
2  & \p Build-in titlepages & little/ can�t include Logo due to license issues\\
%0 & \p altering the rounded theme to have nicer item lables & very little\\
1 & \p improve/add compability to other outer themes than infolines\_hohenheim & little: it seems that e.g. split oder shadow works well with hobete \\
0 & \p transfer to latex3 & non \\
\end{tabularx}\vfill
\tiny
\p = feature to create , \pp =  feature to improve ,  - =  feature to remove , \bee  bug-fixes\\
Importance: 0=non ... 3=very
}




%\mylogo{ \includegraphics[height=5ex,clip=true , trim=8.1cm  0pt 0pt 0pt ]{logo.eps} }
%-- Yes it can be changed within the document




\section{Beispiele}

\sectionpage
\sectionpage[plain=true]

\frame{\frametitle{Tabelle}\small
\begin{tabularx}{.99\textwidth}{XXXXXX}
\toprule
Prozess\-schritt & Tem\-peratur �C & Rel. Luft-feuchte \% & Dreh\-zahl Luftum\-w�lzung & Frischluft-klappen &   Zeit h  \\ \midrule

      1. &         24 &         94 &          1 &          - &      18.00  \\

        2. &         22 &         90 &          1 &          - &      24.00   \\

        3. &         21 &         86 &          2 &          - &      24.00.  \\

        4. &         18 &         84 &          2 &         1/4  &      24.00  \\

        5. &         16 &         80 &          2 &      1/4   &      48.00  \\

        6. &         16 &         76 &          2 &       1/4   &      24.00 \\ \bottomrule
\end{tabularx}
}

\frame{\begin{center}
\begin{tikzpicture}
\begin{axis}[scale only axis, xmin=0.9,xmax=6.1,axis y line*=left, xlabel=Schritt ,ylabel=\textcolor{red}{Temperatur}] \addplot[color=red,mark=x, smooth] coordinates {
(1, 24) (2, 22) (3, 21) (4,18) (5, 16) (6, 16)
};\end{axis}

\begin{axis}[ scale only axis, xmin=0.9,xmax=6.1, axis y line*=right, axis x line=none, ylabel=\textcolor{blue}{rel. Luftfeuchte}]
\addplot[color=blue,mark=square, smooth] coordinates {
(1, 94) (2, 90) (3, 86)  (4, 84) (5, 80) (6, 76)
};
\end{axis} \end{tikzpicture}
\end{center}}
\frame{
\begin{alertblock}{Achtung!}
Guten Tag
\end{alertblock}

\begin{exampleblock}{Beispiel}
Hallo
\end{exampleblock}
\begin{definition}
So ist das.
\end{definition}
\begin{theorem}
So k�nnte es sein
\end{theorem}

}

\frame{
\begin{proof}
Genau so ist es!
\end{proof}


}
\subsection{Listen}
\frame{\frametitle{Beispiele F�r Listen}
\begin{itemize}[<+-| alert@+>]
\item Eins
\item Zwei
\item Drei
\item Vier
\begin{itemize}
\item bla
\end{itemize}

\end{itemize}
\begin{enumerate}
\item Eins
\item Zwei
\end{enumerate}
}
\subsection{s.th. else}
\frame{}




\end{document}
\end{frame}

\begin{frame}[fragile]\frametitle{Die Framelist}
Die Framelist ist eine Liste der Frametitles. Diese wird �hnlich zu toc�s in einer Hilfsdatei mit der Endung \texttt{.frml} abgelegt. Der Befehl \lstinline[style=tex]$ \printframelist $ gibt die Liste aus. Mit der Klassenoption \lstinline[style=tex]$ framelist front$ und \lstinline[style=tex]$ framelist back $ k�nnen die Token eingestellt werden, die vor bzw nach den Eintr�gen erscheinen. Defaultwerte sind \lstinline[style=tex]{ $\cdot$ } und \lstinline[style=tex]$ \\ $ .
\end{frame}


\newcommand{\p}{{\footnotesize +}\ }\newcommand{\pp}{{\footnotesize ++}\ }

\frame{\frametitle{Change-Log}\scriptsize
\begin{tabularx}{\textwidth}{XXX}
Version & Changes & Status-Name\\
0001& \p Colortheme\newline \p Fancy Framtitle\newline \p mylogo-Command  &  beta-hack\\[15mm]
0001a & \bee   improved frametitle\newline \bee correct name for yellow \newline \p better  Logo control \newline \p added \texttt{HohenheimFancyTitle} &\\[15mm]
0001b & \bee   minor bug fixes\newline &\\
0001c & \p sectionpage now with plain option \p Better \texttt{HohenheimFancyTitle} \p altered item labels & \\
0002a & \p added framelist   & beta-hack release candiate (the steped version number indicates that this is a realease version) \\
\end{tabularx}\vfill
\tiny
\p = new feature, \pp = significantly improved feature,  - = removed feature, \bee  bug-fix
}
\frame{\frametitle{Change-Log Vorts.}\scriptsize
\begin{tabularx}{\textwidth}{XXX}
Version & Changes & Status-Name\\
0002b & \p added Support for beamerposter & \\
0003 & \pp tidy up of the code, from now on its a real Package & \\
0003a & \pp growing support for beamerposter & \\
0004 & \pp latex3 plus options & \\
\end{tabularx}\vfill
\tiny
\p = new feature, \pp = significantly improved feature,  - = removed feature, \bee  bug-fix
}





\frame{\frametitle{To Do }\scriptsize

\begin{tabularx}{\textwidth}{XXX}
Importance & Changes & Progress/Probs\\
0& \p Support for beamerarticle &  non \\ 
\rowcolor{black!10} 
2  & \p Build-in titlepages & little/ can�t include Logo due to license issues\\
%0 & \p altering the rounded theme to have nicer item lables & very little\\
1 & \p improve/add compability to other outer themes than infolines\_hohenheim & little: it seems that e.g. split oder shadow works well with hobete \\
0 & \p transfer to latex3 & non \\
\end{tabularx}\vfill
\tiny
\p = feature to create , \pp =  feature to improve ,  - =  feature to remove , \bee  bug-fixes\\
Importance: 0=non ... 3=very
}




%\mylogo{ \includegraphics[height=5ex,clip=true , trim=8.1cm  0pt 0pt 0pt ]{logo.eps} }
%-- Yes it can be changed within the document




\section{Beispiele}

\sectionpage
\sectionpage[plain=true]

\frame{\frametitle{Tabelle}\small
\begin{tabularx}{.99\textwidth}{XXXXXX}
\toprule
Prozess\-schritt & Tem\-peratur �C & Rel. Luft-feuchte \% & Dreh\-zahl Luftum\-w�lzung & Frischluft-klappen &   Zeit h  \\ \midrule

      1. &         24 &         94 &          1 &          - &      18.00  \\

        2. &         22 &         90 &          1 &          - &      24.00   \\

        3. &         21 &         86 &          2 &          - &      24.00.  \\

        4. &         18 &         84 &          2 &         1/4  &      24.00  \\

        5. &         16 &         80 &          2 &      1/4   &      48.00  \\

        6. &         16 &         76 &          2 &       1/4   &      24.00 \\ \bottomrule
\end{tabularx}
}

\frame{\begin{center}
\begin{tikzpicture}
\begin{axis}[scale only axis, xmin=0.9,xmax=6.1,axis y line*=left, xlabel=Schritt ,ylabel=\textcolor{red}{Temperatur}] \addplot[color=red,mark=x, smooth] coordinates {
(1, 24) (2, 22) (3, 21) (4,18) (5, 16) (6, 16)
};\end{axis}

\begin{axis}[ scale only axis, xmin=0.9,xmax=6.1, axis y line*=right, axis x line=none, ylabel=\textcolor{blue}{rel. Luftfeuchte}]
\addplot[color=blue,mark=square, smooth] coordinates {
(1, 94) (2, 90) (3, 86)  (4, 84) (5, 80) (6, 76)
};
\end{axis} \end{tikzpicture}
\end{center}}
\frame{
\begin{alertblock}{Achtung!}
Guten Tag
\end{alertblock}

\begin{exampleblock}{Beispiel}
Hallo
\end{exampleblock}
\begin{definition}
So ist das.
\end{definition}
\begin{theorem}
So k�nnte es sein
\end{theorem}

}

\frame{
\begin{proof}
Genau so ist es!
\end{proof}


}
\subsection{Listen}
\frame{\frametitle{Beispiele F�r Listen}
\begin{itemize}[<+-| alert@+>]
\item Eins
\item Zwei
\item Drei
\item Vier
\begin{itemize}
\item bla
\end{itemize}

\end{itemize}
\begin{enumerate}
\item Eins
\item Zwei
\end{enumerate}
}
\subsection{s.th. else}
\frame{}




\end{document}
\end{frame}

\begin{frame}[fragile]\frametitle{Die Framelist}
Die Framelist ist eine Liste der Frametitles. Diese wird �hnlich zu toc�s in einer Hilfsdatei mit der Endung \texttt{.frml} abgelegt. Der Befehl \lstinline[style=tex]$ \printframelist $ gibt die Liste aus. Mit der Klassenoption \lstinline[style=tex]$ framelist front$ und \lstinline[style=tex]$ framelist back $ k�nnen die Token eingestellt werden, die vor bzw nach den Eintr�gen erscheinen. Defaultwerte sind \lstinline[style=tex]{ $\cdot$ } und \lstinline[style=tex]$ \\ $ .
\end{frame}


\newcommand{\p}{{\footnotesize +}\ }\newcommand{\pp}{{\footnotesize ++}\ }

\frame{\frametitle{Change-Log}\scriptsize
\begin{tabularx}{\textwidth}{XXX}
Version & Changes & Status-Name\\
0001& \p Colortheme\newline \p Fancy Framtitle\newline \p mylogo-Command  &  beta-hack\\[15mm]
0001a & \bee   improved frametitle\newline \bee correct name for yellow \newline \p better  Logo control \newline \p added \texttt{HohenheimFancyTitle} &\\[15mm]
0001b & \bee   minor bug fixes\newline &\\
0001c & \p sectionpage now with plain option \p Better \texttt{HohenheimFancyTitle} \p altered item labels & \\
0002a & \p added framelist   & beta-hack release candiate (the steped version number indicates that this is a realease version) \\
\end{tabularx}\vfill
\tiny
\p = new feature, \pp = significantly improved feature,  - = removed feature, \bee  bug-fix
}
\frame{\frametitle{Change-Log Vorts.}\scriptsize
\begin{tabularx}{\textwidth}{XXX}
Version & Changes & Status-Name\\
0002b & \p added Support for beamerposter & \\
0003 & \pp tidy up of the code, from now on its a real Package & \\
0003a & \pp growing support for beamerposter & \\
0004 & \pp latex3 plus options & \\
\end{tabularx}\vfill
\tiny
\p = new feature, \pp = significantly improved feature,  - = removed feature, \bee  bug-fix
}





\frame{\frametitle{To Do }\scriptsize

\begin{tabularx}{\textwidth}{XXX}
Importance & Changes & Progress/Probs\\
0& \p Support for beamerarticle &  non \\ 
\rowcolor{black!10} 
2  & \p Build-in titlepages & little/ can�t include Logo due to license issues\\
%0 & \p altering the rounded theme to have nicer item lables & very little\\
1 & \p improve/add compability to other outer themes than infolines\_hohenheim & little: it seems that e.g. split oder shadow works well with hobete \\
0 & \p transfer to latex3 & non \\
\end{tabularx}\vfill
\tiny
\p = feature to create , \pp =  feature to improve ,  - =  feature to remove , \bee  bug-fixes\\
Importance: 0=non ... 3=very
}




%\mylogo{ \includegraphics[height=5ex,clip=true , trim=8.1cm  0pt 0pt 0pt ]{logo.eps} }
%-- Yes it can be changed within the document




\section{Beispiele}

\sectionpage
\sectionpage[plain=true]

\frame{\frametitle{Tabelle}\small
\begin{tabularx}{.99\textwidth}{XXXXXX}
\toprule
Prozess\-schritt & Tem\-peratur �C & Rel. Luft-feuchte \% & Dreh\-zahl Luftum\-w�lzung & Frischluft-klappen &   Zeit h  \\ \midrule

      1. &         24 &         94 &          1 &          - &      18.00  \\

        2. &         22 &         90 &          1 &          - &      24.00   \\

        3. &         21 &         86 &          2 &          - &      24.00.  \\

        4. &         18 &         84 &          2 &         1/4  &      24.00  \\

        5. &         16 &         80 &          2 &      1/4   &      48.00  \\

        6. &         16 &         76 &          2 &       1/4   &      24.00 \\ \bottomrule
\end{tabularx}
}

\frame{\begin{center}
\begin{tikzpicture}
\begin{axis}[scale only axis, xmin=0.9,xmax=6.1,axis y line*=left, xlabel=Schritt ,ylabel=\textcolor{red}{Temperatur}] \addplot[color=red,mark=x, smooth] coordinates {
(1, 24) (2, 22) (3, 21) (4,18) (5, 16) (6, 16)
};\end{axis}

\begin{axis}[ scale only axis, xmin=0.9,xmax=6.1, axis y line*=right, axis x line=none, ylabel=\textcolor{blue}{rel. Luftfeuchte}]
\addplot[color=blue,mark=square, smooth] coordinates {
(1, 94) (2, 90) (3, 86)  (4, 84) (5, 80) (6, 76)
};
\end{axis} \end{tikzpicture}
\end{center}}
\frame{
\begin{alertblock}{Achtung!}
Guten Tag
\end{alertblock}

\begin{exampleblock}{Beispiel}
Hallo
\end{exampleblock}
\begin{definition}
So ist das.
\end{definition}
\begin{theorem}
So k�nnte es sein
\end{theorem}

}

\frame{
\begin{proof}
Genau so ist es!
\end{proof}


}
\subsection{Listen}
\frame{\frametitle{Beispiele F�r Listen}
\begin{itemize}[<+-| alert@+>]
\item Eins
\item Zwei
\item Drei
\item Vier
\begin{itemize}
\item bla
\end{itemize}

\end{itemize}
\begin{enumerate}
\item Eins
\item Zwei
\end{enumerate}
}
\subsection{s.th. else}
\frame{}




\end{document}
\end{frame}

\begin{frame}[fragile]\frametitle{Die Framelist}
Die Framelist ist eine Liste der Frametitles. Diese wird �hnlich zu toc�s in einer Hilfsdatei mit der Endung \texttt{.frml} abgelegt. Der Befehl \lstinline[style=tex]$ \printframelist $ gibt die Liste aus. Mit der Klassenoption \lstinline[style=tex]$ framelist front$ und \lstinline[style=tex]$ framelist back $ k�nnen die Token eingestellt werden, die vor bzw nach den Eintr�gen erscheinen. Defaultwerte sind \lstinline[style=tex]{ $\cdot$ } und \lstinline[style=tex]$ \\ $ .
\end{frame}


\newcommand{\p}{{\footnotesize +}\ }\newcommand{\pp}{{\footnotesize ++}\ }

\frame{\frametitle{Change-Log}\scriptsize
\begin{tabularx}{\textwidth}{XXX}
Version & Changes & Status-Name\\
0001& \p Colortheme\newline \p Fancy Framtitle\newline \p mylogo-Command  &  beta-hack\\[15mm]
0001a & \bee   improved frametitle\newline \bee correct name for yellow \newline \p better  Logo control \newline \p added \texttt{HohenheimFancyTitle} &\\[15mm]
0001b & \bee   minor bug fixes\newline &\\
0001c & \p sectionpage now with plain option \p Better \texttt{HohenheimFancyTitle} \p altered item labels & \\
0002a & \p added framelist   & beta-hack release candiate (the steped version number indicates that this is a realease version) \\
\end{tabularx}\vfill
\tiny
\p = new feature, \pp = significantly improved feature,  - = removed feature, \bee  bug-fix
}
\frame{\frametitle{Change-Log Vorts.}\scriptsize
\begin{tabularx}{\textwidth}{XXX}
Version & Changes & Status-Name\\
0002b & \p added Support for beamerposter & \\
0003 & \pp tidy up of the code, from now on its a real Package & \\
0003a & \pp growing support for beamerposter & \\
0004 & \pp latex3 plus options & \\
\end{tabularx}\vfill
\tiny
\p = new feature, \pp = significantly improved feature,  - = removed feature, \bee  bug-fix
}





\frame{\frametitle{To Do }\scriptsize

\begin{tabularx}{\textwidth}{XXX}
Importance & Changes & Progress/Probs\\
0& \p Support for beamerarticle &  non \\ 
\rowcolor{black!10} 
2  & \p Build-in titlepages & little/ can�t include Logo due to license issues\\
%0 & \p altering the rounded theme to have nicer item lables & very little\\
1 & \p improve/add compability to other outer themes than infolines\_hohenheim & little: it seems that e.g. split oder shadow works well with hobete \\
0 & \p transfer to latex3 & non \\
\end{tabularx}\vfill
\tiny
\p = feature to create , \pp =  feature to improve ,  - =  feature to remove , \bee  bug-fixes\\
Importance: 0=non ... 3=very
}




%\mylogo{ \includegraphics[height=5ex,clip=true , trim=8.1cm  0pt 0pt 0pt ]{logo.eps} }
%-- Yes it can be changed within the document




\section{Beispiele}

\sectionpage
\sectionpage[plain=true]

\frame{\frametitle{Tabelle}\small
\begin{tabularx}{.99\textwidth}{XXXXXX}
\toprule
Prozess\-schritt & Tem\-peratur �C & Rel. Luft-feuchte \% & Dreh\-zahl Luftum\-w�lzung & Frischluft-klappen &   Zeit h  \\ \midrule

      1. &         24 &         94 &          1 &          - &      18.00  \\

        2. &         22 &         90 &          1 &          - &      24.00   \\

        3. &         21 &         86 &          2 &          - &      24.00.  \\

        4. &         18 &         84 &          2 &         1/4  &      24.00  \\

        5. &         16 &         80 &          2 &      1/4   &      48.00  \\

        6. &         16 &         76 &          2 &       1/4   &      24.00 \\ \bottomrule
\end{tabularx}
}

\frame{\begin{center}
\begin{tikzpicture}
\begin{axis}[scale only axis, xmin=0.9,xmax=6.1,axis y line*=left, xlabel=Schritt ,ylabel=\textcolor{red}{Temperatur}] \addplot[color=red,mark=x, smooth] coordinates {
(1, 24) (2, 22) (3, 21) (4,18) (5, 16) (6, 16)
};\end{axis}

\begin{axis}[ scale only axis, xmin=0.9,xmax=6.1, axis y line*=right, axis x line=none, ylabel=\textcolor{blue}{rel. Luftfeuchte}]
\addplot[color=blue,mark=square, smooth] coordinates {
(1, 94) (2, 90) (3, 86)  (4, 84) (5, 80) (6, 76)
};
\end{axis} \end{tikzpicture}
\end{center}}
\frame{
\begin{alertblock}{Achtung!}
Guten Tag
\end{alertblock}

\begin{exampleblock}{Beispiel}
Hallo
\end{exampleblock}
\begin{definition}
So ist das.
\end{definition}
\begin{theorem}
So k�nnte es sein
\end{theorem}

}

\frame{
\begin{proof}
Genau so ist es!
\end{proof}


}
\subsection{Listen}
\frame{\frametitle{Beispiele F�r Listen}
\begin{itemize}[<+-| alert@+>]
\item Eins
\item Zwei
\item Drei
\item Vier
\begin{itemize}
\item bla
\end{itemize}

\end{itemize}
\begin{enumerate}
\item Eins
\item Zwei
\end{enumerate}
}
\subsection{s.th. else}
\frame{}




\end{document}